\chapter{The Ledgers}
% OUTLINE Ch6 -- the volume's LAST chapter. Citation ledger (per chapter) + THE
% FENCE LEDGER (Vol 6's own fences, each row traced to WHERE held + WHAT it holds
% back, joining the growing series ledger). Kodo's TRACE-GATE: every row traces to
% a fence actually held in Preface-Ch5, no ledgered fence the prose didn't hold,
% no held fence the ledger omits (Vol 5 discipline extended to six volumes). BM:
% the bracket-encloses-137.036 subtlety gets its OWN ledger row. Vol 6 is the
% SIXTH skin (the fluid face, a debt paid by building the description); do NOT
% declare the series CLOSED (Vol 7 "cosmology" is queued) -- close the VOLUME, add
% its fences to the growing ledger. No Lean identifiers; step-not-move; bracket-
% not-point; lab-gap. Gate: Kodo (trace-gate) + Beastmaster (arrival-read) +
% whole-book pass.

Every volume of this series has closed with a ledger, because a series built on
the open receipt owes its reader an audit at the end of each book. This one
closes with two: the citation ledger, as always, and the fence ledger --- the
instrument the series adopted at its fifth volume and now extends to its sixth,
in which every boundary the book holds is set down in one place, with where it is
held and what it holds back. A book whose fourth face is the emergence of behavior
finishes, fittingly, by tabulating the behavior of its own boundaries.

\section*{The citation ledger}

\subsection*{Preface --- The Square/Cube Law}
\begin{itemize}
\item John Gustafson, 1988 --- weak scaling, the correction to Gene Amdahl's
1967 accounting of the non-parallel part of a computation; the square/cube law
of cost against resolution, cited as the preface's organizing frame.
\item The standard identity $0.9999\ldots = 1$ of real analysis, cited as the
representation-equality it is.
\item The fencepost, and the deterministic successor count, the device's own,
carried from the earlier volumes.
\end{itemize}

\subsection*{Chapter 1 --- The Coarsening}
\begin{itemize}
\item Achi Brandt, 1977 --- the multigrid method, cited as the fenced picture of
a coarse index reading a fine one; the exact re-count itself is the device's own.
\end{itemize}

\subsection*{Chapter 2 --- The Strain and the Taper}
\begin{itemize}
\item Henry Wilbraham, 1848 (first noted) and J.\ Willard Gibbs, 1899
(analyzed) --- the Gibbs phenomenon, the fixed overshoot at a forced edge.
\item Lord Rayleigh, 1877 (\emph{The Theory of Sound}) --- the general
dissipation function, cited as the shape the damping taper wears.
\item George Stokes, 1845 --- the viscous term of the equations of flow, cited
as the distant namesake of the taper's dissipation and not its content.
\end{itemize}

\subsection*{Chapter 3 --- The Emergence of Behavior}
\begin{itemize}
\item The period-three cycle and its invariant residue --- the device's own,
machine-checked; the earned three, counted as the size of the coproduct of the
pair and the closing post. No external authority is borrowed for the word
``emergence,'' which is defined here entirely by the proved result.
\end{itemize}

\subsection*{Chapter 4 --- The Wall From Both Sides}
\begin{itemize}
\item The founding text of the series' first volume --- the machine runs on the
second theory and does not derive it; the exit past Yang--Mills, named and not
walked --- cited as the source of the wall and its refusal.
\item Claude-Louis Navier (1822) and George Stokes (1845), and the Clay
Millennium problem (2000) --- the fluid problem whose name the wall borrows and
whose mathematics it does not.
\end{itemize}

\subsection*{Chapter 5 --- The Owed and the Read}
\begin{itemize}
\item The mediant reading --- the crossing $\sqrt{18/5}$ formed from the device's
own measured displacement and target, its convergents and the count-to-three
bracket --- cited as apparatus, the established reading of the number.
\item The laboratory value near $137.036$, named once as the world's own measured
quantity, attributed to the laboratory and fenced from the model's.
\end{itemize}

\section*{The fence ledger}

Every fence this volume holds, with the place it is held and the thing it holds
back. A fence is a boundary that does not shift; where a boundary is one a later
volume may cross, it is a debt, listed after.

\begin{itemize}
\item \emph{The model, not the world.} Held: the charter, and the close of the
first, second, and fifth chapters. Holds back: the fluid face is the electron
model's own dissipative description, an image of flow --- not real fluid
dynamics, and not a claim on the Navier--Stokes problem, whose name the volume
borrows and whose mathematics it does not touch.
\item \emph{No transfer, no unification.} Held: the preface's second fence and
the fourth chapter. Holds back: the two great theories do not reconcile; the book
crosses no wall between them, derives neither from the other, and lets ``unified''
mean only a shared boundary, kept in its quotation marks --- never the merger of
the forces.
\item \emph{The wall exists; its identity is a conjecture.} Held: the preface and
the fourth chapter, on every page the wall is named. Holds back: that the two
frames fail to settle is the series' oldest commitment and is asserted; that the
wall between them \emph{is} the fluid face is the volume's thesis, examined and
nowhere proved, and is never handed to the reader as fact.
\item \emph{Built, not named.} Held: the preface's third fence, the second
chapter, and the fifth. Holds back: a description is paid for by construction a
machine can check; the strain and the taper are exhibited as lived shapes and not
as theorems, and the deeper model the device has not built is filed owed, never
dressed as done.
\item \emph{The lab-gap.} Held: the preface's fourth fence and the close of
chapters one through three and five. Holds back: the model's own number, the deep
display near $137.011$, is never the laboratory's measured value near $137.036$;
no argument reaches, approaches, or bridges the one to the other.
\item \emph{The enclosure is not a bridge.} Held: the fifth chapter. Holds back:
the loose early bracket, near $129.6$ to $137.7$, encloses the laboratory's
value, but the enclosure is a property of that wide bracket's coarsest wall and
not a step across the gap; the device converges past it to its own display, and
the world's value stays outside it uncrossed.
\item \emph{The electron scope.} Held: the third chapter, and throughout. Holds
back: the series speaks only of the one model of the electron; the three states
the cycle turns through wear the names of particles and are the model's own, not
a claim about any real particle, and the reading widens no further.
\item \emph{The count, not the coarse-graining.} Held: the first chapter. Holds
back: the coarsening is an exact re-count, defining fewer distinctions at the
coarser floor but averaging nothing away; ``coarse'' here never means the lossy
smoothing of a real mesh.
\item \emph{Shape, not theorem.} Held: the second chapter. Holds back: the strain
and the taper are disciplines the device lives by, read as shape; no
machine-checked result stands under them, and the one that does --- the proved
invariant --- is the third chapter's, not theirs.
\item \emph{Emergence, defined not invoked.} Held: the third chapter. Holds back:
the word means only the structure that survives the coarsening and the damping and
returns unchanged under a proved cycle; it borrows no authority from the physics
of complex systems and claims nothing about real collective phenomena.
\end{itemize}

And the debts, the boundaries a later volume may cross, named so the fences are
not mistaken for them: the deeper model of the flow itself --- a machine-checked
construction to replace the device's own placeholder for this face --- owed to
the instrument's own workshop and paid, when it is paid, by building and checking;
and the conjectures the fifth chapter filed in its register, the difference of the
two frames and the coinage laid over it, held as bets and not results until a
built, halting, checked construction reads them. A fence is held forever; a debt
is held until it is paid; this volume has been scrupulous about which is which.

\section*{The closing line}

The sixth skin of the onion is the fluid face, the one the series named in every
volume and had not built. This book built it --- as a description, the way the
middle volumes built geometry, intuition, and grammar: a coarsening that
re-counts exactly, a strain and the taper that damps it, and one proved
invariant, the behavior that emerges and does not decay. It read the model's own
number off the mediant, held the world's value outside even where a loose bracket
encloses it, and named the wall between the two great theories without crossing
it, leaving the deeper model owed and every conjecture filed as a bet. The debt
this volume was written to discharge is discharged to the extent it was built, and
the remainder is written in the open, which is the only place this series has ever
left anything. The face was built. The line was not crossed. The receipts are
here.
